\section{Introducción: de la generación incondicional a la condicional}

Esta clase es la sexta del curso KAIST CS492D (Modelos de difusión y sus aplicaciones), impartida por el profesor Minhyuk Sung. Las clases anteriores ya presentaron los principios fundamentales de modelos de difusión como DDPM y DDIM, y los estudiantes tienen capacidad para implementar modelos básicos de difusión. En esta clase "cambiaremos de marcha" para discutir cómo introducir \textbf{información condicional} (etiquetas de clase, descripciones de texto, imágenes, etc.) en aplicaciones prácticas y cómo realizar un \textbf{fine-tuning eficiente} de modelos de difusión preentrenados a gran escala.

\begin{knowledgebox}{Repaso de conocimientos previos}
En el marco DDPM/DDIM, existe una relación determinista entre la salida de la red de predicción de ruido $\epsilon_\theta(x_t, t)$ y la función score:
$$
\nabla_{x_t} \log p(x_t) \;\propto\; -\frac{1}{\sqrt{1-\bar\alpha_t}}\,\epsilon_\theta(x_t, t)
$$
donde $\bar\alpha_t$ es el producto acumulado de la programación de ruido. Esta relación constituye la base matemática de todos los métodos de generación condicional: guiar las trayectorias de generación mediante la modificación de la función score.
\end{knowledgebox}

Esta clase cubre los siguientes temas centrales:
\begin{enumerate}
    \item \textbf{Classifier Guidance}: guía condicional mediante un clasificador externo
    \item \textbf{Classifier-Free Guidance (CFG)}: guía condicional sin clasificador, estándar industrial
    \item \textbf{DDIM Inversion}: inversión determinista y edición de imágenes
    \item \textbf{Latent Diffusion Models}: reducción del costo computacional ejecutando la difusión en el espacio latente
    \item \textbf{ControlNet}: control condicional mediante fine-tuning eficiente
    \item \textbf{LoRA}: ajuste de rango bajo para personalización en modelos de difusión
\end{enumerate}

\newpage

